\documentclass[sigconf,authordraft]{acmart}
\usepackage{listings}
\usepackage{amsmath}
\begin{document}

\title{Omelets Need Onions \\ \large E-graphs Modulo Theories via Bottom Up E-Matching}
% Bottom Up E-matching for E-graphs Modulo Theories
% EGraph Modulo THoeires and Bottom Up E-Matching
% E-Graphs Modulo Theories by Bottom Up E-Matching

% bottom up is really a key point in everything.

\author{Philip Zucker}
\email{philzuckerblog@gmail.com}


\begin{abstract}
    E-graphs\cite{eggpaper} are a data structure for generic equational reasoning and optimization over ground terms. One of
     the benefits of e-graph rewriting is that it can declaratively handle useful but difficult to orient identities like 
     associativity and commutativity (AC) in a generic way. However, using these generic mechanisms
      is more computationally expensive than bespoke routines. There exist specialized efficient
       algorithms and data structures for terms containing sets, multisets, linear expressions, polynomials, and binders.
        A natural question arises: How can one combine the generic capabilities of e-graph rewriting with these specialized theories?
        This talk discusses a pragmatic approach to this e-graphs modulo theories (EMT) question using two key ideas: bottom up e-matching and semantic e-ids.

% One feels there must be bespoke reoutines that would be better suited for common problems.
% Some theories feel nearly structural.

\end{abstract}


\keywords{e-graphs, e-matching, term rewriting, equality saturation, SMT, union find, completion, theory solvers, Gr\"{o}bner basis}


\maketitle

%\section{Introduction}

% Consider equality saturating the AC axioms $\forall x y, x + y = y + x$ and $\forall x y z, (x + y) + z = x + (y + z)$  to the expression $x_1 + (x_2 + (x_3 + (... + x_N)))$. Upon saturation e-graph will 
% contain $2^N - 1$ distinct e-classes, one for each non-empty subset of $\{x_1, x_2, ... \}$. The e-graph will also contain 
% $3^N - 2^{N+1} +1$ distinct enodes. Clearly the e-graph is exponentially exploding in terms of the size of these expressions, but for trivial reasons. \cite{souffleegraph}

% This explosion is in one sense unavoidable. All the e-classes are semantically distinct entities we may actually want to examine in the general case.
% What we can do however is produce a pruned subset on demand rather than materialize ahead of time. 
% % We can also perhaps reduce the number of enodes by baking in.
% The e-graph automatically prunes considered terms to only those that are subterms of a member of the equivalence class of the goal term.


% The beauty of greedy rewriting is that it has both an operational and declarative semantics. Likewise can be said for logic programming.
% These systems tweak the knob pretty far to operational reading and it is on programmer to avoid blind alleys and loops.

% If one only wants a system that is complete but brute force, there are many ways to do this and available ATPs even are generically perfomant.
% What we want is a system with simple operational mechanism that can be steerable

% E-graph rewriting hits a sweet spot.


% We really may want to examine all of these possibilities in. 


% What do I mean by "theories"

% ground terms are easier to deal with than terms with variables in them. Patterning matching is solving an equation bwetween a varterm and a term.
% The main reason to prefer pattern matching is that it is more tractable in many cases and conceptually simpler.

\section{Terms Modulo Theories}

It is instructive to dial back from the complexity of e-graphs to see how one can make sense of the "baking in" of theories for ordinary tree-like terms.

The most basic form of a term is a function symbol and an ordered list of children
\begin{lstlisting}[language=Caml]
type term = Fn of string * term list
\end{lstlisting}

We can extend this to a more general form of term with specialized containers of children

% cut this? It's out of place now
\begin{lstlisting}[language=Caml]
type term = 
  | Fn of string * term list
  | MS of string * term multiset
  | Set of string * term set
  | Poly of string * term poly
  | LitInt of int
\end{lstlisting}

% type container = List of term | list MS of list | Set | Poly
% type term = Fn of string * container

% term = Default of term list | A of term list | AC of multiset | ACI of set | Ring of term poly | Value of int

%or
% type node_type = MS | List | Set   %| Poly | 
% type term = Fn of string * node_type * term list

% \begin{itemize}
% \item structural equality
% \item $\alpha$ equality
% \item $\beta_0$ equality
% \item $\beta$ equality
% \item observational equality
% \item equality modulo an equational theory
% \item equality modulo a rewriting theory
% \end{itemize}

The specialized data structures like set and multiset can be implemented using regular lists. Structurally canonical list representatives can be found by using sorting and deduplication \cite{hashingmodtheories}. A more performant thing to do can be to use Patricia tries, which are also structurally canonical but support faster set and map operations.

One can also use smart constructors \cite{smartconstructor} to apply built in rewrite rules immediately upon construction. They check all rules which have the smart constructor symbol at the head of the pattern. If the rewrite rules are convergent, one can maintain the invariant the terms are only created in canonical form. %This can be mixed and matched with the previous approach, as the container structures can be encoded into the terms.

If the terms are stored in a structurally canonical way, one can use regular hash consing \cite{hashconsing} to intern them. This supports fast equality checks modulo the theories.

This shows that equality checking and interning terms modulo sufficiently nice theories is not hard, but what about pattern matching?

%\cite{maziarz2021hashingmoduloalphaequivalence}

%\subsection{Hash Consing Modulo Theories}

%  A key question about any data structure is what operations do you want support and what has to be fast.
%  We may want to
%  \begin{itemize}
%  \item Compare terms for different notions of equality
%  \item match specific terms against patterns
%  \item Set and Dictionaries of terms
%  \item seek terms in a termbank to with specific properties
 
% % \item Rewrite terms.
% % \item Construct terms
% % \item Canonize terms
% % \item Unify
% % \item Subsumption
% % \item Subterms
% % \item Immediate Children
 
% \end{itemize}

% An e-graph can be described as a hash cons smashed together with a union find, so it is interesting to consider what Hash Consing Modulo Theories might require

% It can be useful to have interned and non-interned terms as distinct data types \cite{kontroliefficient}. Smart constructors may rewrite the term through multiple temporary intermediaries.

% There can be terms that are usefully persisted the entirety of the program, others that may exists an an arena, or further others that may be uniquely owned by their parent.


% Modular hash consing

% Conversion to normal forms can embed into "boring" term. Sort. Particia trees

% Smart constructors







% put bottom up matching here?



\subsection{What patterns are possible?}

Unification and pattern matching are a logical and syntactic slant on the notion of equation solving \cite{traatbook} \cite{BAADER2001445}. The problems are formulated as goal equations containing variables $X, Y,...$ to be solved, and the solutions are substitutions, which are mappings from variables to terms. They operate as backwards reasoning method that breaks down goal equation sets until it is refined into a substitution.
% Pattern matching is a subclass of unification where one side of the equation has nothing to solve for. 

E-matching in the e-graphs community may have a connotation of matching a pattern against an e-graph, but in the term rewriting community it has the connotation of solving a pattern matching problem modulo an equational theory E consisting of a set of universally quantified equality axioms. These are the same notion. The e-graph version of e-matching is related to e-matching a term modulo the theory of ground equations that the e-graph represents.

% Gauss elimination has isloate and eliminate vs add tgoether equations.

Equational theories fall into different classes of how many solutions they may allow for their E-unification or E-matching problems.
Some theories will only have a finite set of solutions, some a countably infinite set, and others an uncountably infinite set.

First order syntactic unification is equation solving for a very rigid, simple structural theory of ordered terms. It is unusual in the spectrum of pattern matchers in that it either succeeds once or fails. For example, the goal $cons(X,Y) =^? cons(zero,nil) $ is solved by the substitution $\{X \mapsto zero, Y \mapsto nil\}$ and the goal $cons(X,X) =^? cons(zero,nil)$ has no solution.

%  \lstinline{val pmatch : pat -> term -> subst option}.

Another case may be that a finite number of solutions must be returned. This is the case for some structural theories like AC symbols, sets and multisets. The goal $\{X , Y\} =^? \{1, 2\}$ is solved by $\{X \mapsto 1, Y \mapsto 2\}$ and $\{X \mapsto 2, Y \mapsto 1\}$. 
%The goal $X \cup Y =^? \{1, 2, 3\}$ is solved by $\{X \mapsto \{1, 2\}, Y \mapsto \{3\}\}$ and $\{X \mapsto \{1, 3\}, Y \mapsto \{2\}\}$ and $\{X \mapsto \{2, 3\}, Y \mapsto \{1\}\}$.

%\lstinline{val pmatch : pat -> term -> subst list}

%\item The multiset equation $\{\{x, y\}\} =^? \{\{1, 2, 2\}\}$ fails

For other theories, an enumerably infinite number of solutions is possible. $X + Y =^? 42$ considered over the integers has solutions $\{ \{x \mapsto n, Y \mapsto 42 - n  \} | n \in \mathbb{Z} \}$.  $X \cap Y =^? \{1,2,3\}$ has solutions $\{ {X \mapsto A, Y \mapsto B} | \{1,2,3\} \subseteq A,B \} $ .

%Pattern matching has a signature \lstinline{val pmatch : pat -> term -> subst stream}.



% \item $x \times y =^? $

Finally, even this may not be possible for some theories. Consider solving $X + Y =^? \pi$ over the $\mathbb{R}$. The only remaining possibility is to instead return a constraint, which more or less just returns the pattern.

In short, e-matching modulo theories on terms runs the spectrum from expensive to hopeless.

%\lstinline{val pmatch : pat -> term -> constraint}


% could put references to E/R stuff here

% What solutions does $x_? + ?y =^? 42$.  % $\{ \{x \mapsto 42 - n, y \mapsto 42 + n\} \mid n \in \bN \}$ 
% This cannot be finitely solved as substitutions.

% $?x + ?y = 42, ?x = 16$ this is uniquely solvable.


%Patterns may look like

% \begin{lstlisting}[language=Caml]
% type pat = 
%   | PFn of string * pat list
%   | PVar of string
%   | PMS1 of string * pat multiset
%   | PMSUnion
%   | PSet of string * pat set
%   | SetUnion string * pat multiset
%   | SetIntersection of string * pat multiset
%   | PPoly of string * pat poly
%   | PLitInt of int


% let pmatch (p : pat) (t : term) : solution = 


% \end{lstlisting}



\section {Bottom Up E-Matching Plays Nicer With Theories}

One way to win is to refuse the play the game.

In the statement of the pattern matching problem, one is given a set of equations $\{p_1 =^? t_t, p_2 =^? t_2, ...\}$ to be solved or a pattern to be matched against a particular term.

A different formulation of the pattern matching problem that solves the issues of top down matching is the bottom up matching procedure. Instead of being given a term and a pattern to match it to, one is instead given a term bank $T$ which scanned to attempt to fill the pattern variables to find a match inside the same termbank $T$. \cite{bottomup} This is a brutally simple method.

It is important to note that while an e-graph can be seen as a compressed representation of a set of ground equations $E$, really it is more like the pair of a set of equations and a termbank $(E, T)$. For example, we may consider the set of equations  $\{foo(a) = b, c = c\}$ to be the same as the set $\{foo(a) = b\}$ since the equation $c = c$ has no interesting content. Nevertheless, if we assert these two sets into the e-graph, we get different e-graphs, one which has the an e-class and e-node for $c$ in its termbank $T$ and one which does not.

This python code shows a shallow embedding of bottom up e-matching. It simple scans choices of variables in the form of for loops. Compared to a virtual machine implementation like in top-down e-matching, this is so simple it does not need a library. It can be an idiom.

\begin{lstlisting}[language=Python]
  # A shallow bottom up embedding of the rewrite
  # foo(bar(X), Y) -> biz(X)
  for X in eclasses:
      for Y in eclasses:
          try:
              # possibly failing lookup in hashcons
              lhs = foo[bar[X], Y]  
              # construct rhs
              rhs = biz(X)          
              # set equal
              egraph.union(lhs,rhs) 
          except e:
              pass
  \end{lstlisting}

%However, in equality saturation one is matching against an entire term bank in the form of the e-graph. This problem can be reduced to the previous problem by a scan of all terms, but this is not necessarily the only method. 

%Vectorized

%Given a termbank T and a pattern p, find substitutions $\{v \mapsto t | t \in T\}$ such that  

For rough performance comparison of the algorithms, consider an e-graph with $E$ e-classes and $N$ e-nodes. On average, each e-class has $N/E$ e-nodes.

Top-down e-matching pays a factor $E$ for a scan of all e-classes to start to pick a root e-class to pattern match over. Every time top-down e-matching expands an e-class into it's e-nodes, it adds a multiplicative factor due to this searching. Because of this, top-down e-matching is exponential in cost in terms of depth of the pattern $d$. The cost of top down e-matching is something like $O(E(\frac{N}{E})^d)$.

Bottom-up e-matching pays for a scan for the number of variables $V$ in the pattern. The constructing the instantiated pattern is just a lookup and pays the cost of a hash table or indexing tree, which is typically something like $O(\ln(N))$. The cost is therefore something more akin to $O(E^V d \ln(N) )$.

For different patterns or termbanks, top-down or bottom-up may be more performant.

% consider maybe using this pattern on a given egraph. foo( foo( foo( X )))
A bad pattern for top down e-matching would be a deep one with few variables. For example, 

$$foo(foo(foo(foo(foo(foo(X))))))$$.

A bad pattern for bottom up e-matching would be a shallow one with many variables. For example,

$$foo(bar(X), bar(Y), bar(Z), bar(W), bar(V))$$

A reason that top down matching may feel natural is that makes sense for pattern matching against a pointer chasing implementation of single term. This is the typical situation in many programming language like contexts.

Another reason bottom up e-matching may not feel natural is that the most commonly considered termbank is one consisting of of ground syntactic terms with no theories. In this case $\frac{N}{E} = 1$ and the exponential character of top down matching is gone. It is a dominant method. I believe that people do not necessarily use simple syntactic terms because they do not want baked in theories. They do it because baking in theories is complex and expensive. Bottom-up e-matching ameliorates these issues.

Bottom up e-matching can be seen as a particular query plan for relational e-matching using a generic join \cite{wang2023freejoinunifyingworstcase}. Likewise Top down e-matching can be seen as the same. 

The positive perspective of relation e-matching \cite{zhang2022relationalematching}  is that it gives one freedom to find a good efficient query plan. A negative perspective of relational e-matching is that it doesn't commit to a choice, leaving behind a vague gesturing to the new problem of query optimization which is yet another subsystem to design and implement. %asbtains

%concrete numbers for number of operations on an actual query
% better yet wqould be benchhmrks also, but that's not happening this week.

Bottom up-matching has the killer feature is that it proceeds in the same direction as adding terms into the egraph. Therefore it uses the same indexing structures and canonization that already exist. Bottom up e-matching grounds pattern variables as quickly as possible, and ground terms are the e-graph's bread and butter. Ground is good.

Top-down e-matching requires the extra ability to expand from e-classes to e-nodes and from e-nodes to their children. This is as we have seen, extra challenging or even impossible in the face of theories. Top-down keeps the pattern in an ungrounded state as long as possible. Unground is bad.

Bottom up e-matching can emulate the results of top-down e-matching. If it is possible to enumerate the subterms of your target term, you can fill a term bank with these subterms. Bottom up e-matching will then find a superset of the same matches top-down will. In the presence of some theories enumerating subterms is not possible or meaningful and top-down e-matching will not work at all. You can still bottom up match over a term databank. In this sense, bottom up e-matching is more powerful than top-down e-matching.



% \begin{lstlisting}[language=python]
% for X in eclasses:
%     try:
%         t1 = bar[X] # This can be lifted into higher loop
%         for Y in eclasses:
%             try:
%                 lhs = foo[t1, Y]  # possibly failing lookup in hashcons
%                 rhs = biz(X)          # construct rhs
%                 egraph.union(lhs,rhs) # set equal
%             except:
%                 pass
%     except:
%         pass
% \end{lstlisting}

The maximally naive bottom up e-matching that scans over all e-classes can be improved by filtering on the intersection of all the tables where the variable occurs and by pushing intermediate constructions upwards in the loop ordering. 

\begin{lstlisting}[language=python]
for X in set(bar.keys()).intersection(biz.keys()):
  try:
    t1 = bar[X]
    for Y in set([k[1] for k in foo.keys()]):
      try:
          lhs = foo[t1, Y]
          rhs = biz(X)
          egraph.union(lhs,rhs)
\end{lstlisting}

% the distinction between encapsulate and non-enscpuslated thoery expressions.
One can also filter by modelling using multisorted logic. One only needs to scan over the e-ids corresponding to the sort of te variable. This can be achieved by having separate union finds per sort. Which brings us into our next topic.

\section {Generalized Union Finds and Semantic E-ids}
%E-graphs systems are the interplay of two entities, e-nodes and e-classes. A natural approach to generalize e-graphs is to consider generalizing one or both.

After bottom up e-matching, the second key ingredient for e-graphs modulo theories is to generalize the e-id and the union find. E-ids are generalized to structured e-ids / semantic e-ids / notions of values. The union find becomes a theory specific canonizer that supports equality assertions.

A first way to see this might make sense is to note that a simple e-graph can be seen as a uninterpreted partial-function minimal model $\mathcal{M}$ of the equations and term existence assertions that produced it. The collection of e-ids $\mathcal[V]$ corresponds to the set of values in the model and the e-node tables correspond to the data of the model of the partial functions $\lBrack f \rBrack$. One can ask for interpretation onto domains other than an uninterpreted one where the values have no structure except equality. One can consider semantic or structured domains like polynomials, linear expressions, terms, constructors, sets, multisets and so on. These corresponds to structured e-ids. Multiple kinds of structured e-ids can be most simply supported by using multi-sorted logic. Expressions of different sorts are restricted by being equated by a static type check and the role of mediating between these sorts is performed by the e-nodes and congruence closure through e-nodes.

Another point is to note that even in hash consing for terms modulo theories, there is utility in having both interned and non interned forms of terms. A smart constructor should not bother interning the intermediate expressions until they have been canonized. The e-nodes are the interned forms and the structured e-ids are the non-interned forms. The structured e-ids are in a sense kept in extracted non-interned term-like form for easy deep inspection and theory specific destructive manipulation. While fully interning makes for the simplest fastest equality checks, this is a knob that can be turned between complete structural traversal and a simple id comparison. 

% maybe a nice little picture would be nice here
% in smt and egraph and sat solver mediate all the theories.
% the architecture I'm suggesting is enodes <-> theory1 <-> enodes . Enodes kind of look like propagators networkx.


%The purpose of the union find in an e-graph is canonize e-class ids. It does this by providing the following api/interface.

% take a look at ocaml Fillatire union find api https://ocaml.org/p/unionFind/20220122/doc/UnionFind/Make/index.html

% compare and constrast with smt theory interface. No push pop. No propagate?
% when one puts it like that, it's almost a boring observation

% maybe I should just use python. I haven't prototype this in ocaml
% how do you get theory specific operations? 

% \begin{lstlisting}[language=Caml]
% module type UnionFind
%   type t
%   type eid
%   val create : unit -> t
%   val eq     : eid -> eid -> bool
%   val fresh  : t -> eid                    (* makeset *)
%   val canon  : t -> eid -> eid             (* find    *)
%   val assert_eq : t -> eid -> eid -> unit  (* union   *)
%   val rebuild   : t -> unit
  
%   val hash : t -> eid -> int
%   val compare : t -> eid -> eid -> int
%   val explain : t -> eid -> eid -> explanation

% \end{lstlisting}



% propagate : t -> eid -> eid -> eid list   eids that may have changed?
% propagate ~ rebuild ?

% subterms : t -> eid -> eid list   sometimes is available like for AC or A.
% terms : t -> eid list             if maintains own termbank

% add_term : t -> eid -> unit  % only makes sense if eid was constructed using
% remove_term : t -> eid -> unit % can remove terms from termbank. Makes them unavaiable for ematching, but does not remove their equations

Another way to see that this makes sense is to ask "what interface does the union find present?". The internals of the union find and the integer representation of the e-ids are irrelevant to how they are used in the e-graph. Really one only needs the abstract interface

% module type UnionFind
\begin{lstlisting}[language=Caml]
type t
type eid
val create : unit -> t
val eq     : t -> eid -> eid -> bool
val fresh  : t -> eid (* makeset *)
val canon  : t -> eid -> eid (* find *)
val assert_eq : t -> eid -> eid -> unit (* union *)
\end{lstlisting}

This API is largely the same as that required by the theory of an SMT solver minus the support for backtracking. It is particularly close to the interface of Shostak style theory combination. This supports the the schematic equation $EMT = SMT - SAT$.

\subsection{Union Finds solve Ground Atomic Equations}

The union find can be seen as a normal form or canonizer of a system of atomic equations. The Knuth Bendix completion of a ground atomic equational theory is guaranteed to terminate and the resulting rewrite system can be interpreted as a union find \cite{groundegraph2024}. 

Consider for example the atomic ground equation system 

\begin{equation}
\begin{split}
e_1 = e_2 \\
e_2 = e_3  \\
e_2 = e_1  \\
e_3 = e_1  \\
e_4 = e_5  \\
\end{split}
\end{equation}

Running Knuth Bendix completion with an atomic term ordering $$e_i < e_j \triangleq i < j$$ results in an atomic ground rewrite system

\begin{equation}
\begin{split}
e_2 \rightarrow e_1 \\
e_3 \rightarrow e_1 \\
e_5 \rightarrow e_4
\end{split}
\end{equation}

The $\rightarrow$ can be represented by a parent pointer in a union find.

% Atomic eclass ids have a minimal amount of structure. THere are other more theory specific solvers that support

% \begin{itemize}
% \item eq
% \item canon / find
% \item assertion of equalities union / assert\_eq
% \end{itemize}

However, atomic ground theories are not the only ones that are guaranteed to be completable to canonizers. There is a spectrum of theory specific "union finds" that support canonization and ground equality assertion \cite{MARCHE1996253}. Many of them have the flavor of a specialized theory specific ground Knuth Bendix completion.

\subsection{Example Theories}

\begin{itemize}

\item  Linear equations can be put into row echelon form. The row echelon form can be used to canonically quotient a vector space by a subspace. a set of linear equations can be put into the matrix form $Ax=0$. Each equation of the above example system can be encoded as a row of a matrix. 

\[
\begin{bmatrix}
1 & - 1 & 0 & 0 & 0 \\
0 & 1 & -1 & 0 & 0 \\
-1 & 1 & 0 & 0 & 0 \\
-1 & 0 & 1 & 0 & 0 \\
0 & 0 & 0 & 1 & -1 \\
\end{bmatrix}
\begin{bmatrix}
   e_1 \\
   e_2 \\
   e_3 \\
   e_4 \\
   e_5 \\
\end{bmatrix}
 = 0
\]

The row echelon form corresponds to a ground rewrite system. It can be used starting from the bottom to canonically normalize $\sum_i a_i e_i $ expressions where $a_i \in \mathbb{Q}$. For example $e_1 + 3e_3 + 5e_5 \rightarrow e_1 + e_3 + 5e_4 \rightarrow e_1 + 3e_2 + 5 e_4 \rightarrow 4e_1 + 5e_4$. The method however works just as well with asserted equations like $4e_1 + 5e_4 = 13e_2 + 2e_3$, extending the capabilities of the union find to bake in linear expressions.

\[
\begin{bmatrix}
1 & - 1 & 0 & 0 & 0 \\
0 & 1 & -1 & 0 & 0 \\
0 & 0 & 0 & 1 & -1 \\
\end{bmatrix}
\begin{bmatrix}
   e_1 \\
   e_2 \\
   e_3 \\
   e_4 \\
   e_5 \\
\end{bmatrix}
 = 0
\]

% ⎡1  -1  0   0  0 ⎤
% ⎢                ⎥
% ⎢0  1   -1  0  0 ⎥
% ⎢                ⎥
% ⎢0  1   -1  0  0 ⎥
% ⎢                ⎥
% ⎣0  0   0   1  -1⎦

% ⎡1  0  -1  0  0 ⎤
% ⎢               ⎥
% ⎢0  1  -1  0  0 ⎥
% ⎢               ⎥
% ⎢0  0  0   1  -1⎥
% ⎢               ⎥
% ⎣0  0  0   0  0 ⎦

The semantic domain is a vector space of uninterpreted e-ids quotiented by the subspace spanned by the asserted linear equations. $\mathcal{V} = \mathrm{span}_\mathbb{Q} \{ e_1, e_2, e_3, ... \} \Big/ V_E$ 


\item Gr\"{o}bner bases are the generalization of row echelon form the multivariate polynomial equation systems. This is particularly useful because much of trigonometry and Euclidean geometry can be put into a form that is essentially polynomials for example $\sin(x)**2 + \cos(x)**2 = 1$.   The semantic values are the quotient ring of polynomials modulo the ideal generated by the polynomial equations. $\mathcal{V} = \mathbb{Q}[e_1,e_2,..] \Big/ I_E$
%The cross, geometric, and dot product of vectors are all algebraic constructions. Parallelism, circles, spheres. It is the connection between Radian angles and lengths which requires transcendental functions or solving differential equations or integrals. %Wildberger

\item Ground multiset equations are completeable $ \{\} a,b,b \}\} = \{c,d\}$. The structured e-ids are multisets. This gives each sort one special associative and commutative operation. \cite{acegraph} This can be seen as a special case of Gro\"{o}bner basis calculation as each monomial $x^n y^m z^k$ can be seen as a multiset of it's variables with the multiplicity given by the exponent. Indeed, the termination of Buchberger's algorithm is guaranteed by multiset orderings on the produced polynomials. It can also be seen as integer lattice rewriting via Hilbert or Graver bases and has connections to the theory of integer programming \cite{sturmfels1996grobner}

\item Ground terms equations. In an amusing self reflection, the theory solver could itself be an e-graph, which is a canonizer for ground term equations. It is unclear why this would be useful. One could separate functions symbols to those considered outer and inner. This would enable for example, two e-graph implementations to interact with each other.

\item Atomic Completion modulo a group action. A union find can be built that carries group elements on it's edges \cite{Frühwirth_2009} \cite{lemerre2024labeled} \cite{groupoidunion}. The composition law of the group enables composition while chasing up the parent pointer of the union find, and group inversion enables turning comparison to the root back into comparison with a leaf of the union find. In this manner it can be used to store an equivalence relation modulo a group action. For example the group of integer addition `x + 6 = y - 3` while subsumed by the Row echelon method can be more simply and efficiently implemented using this technique. Permutation groups are also very useful in the context of nominal rewriting and slotted egraphs.
% Another way of looking at it is atomic ground completion modulo a group action.
% Kmett, CHR book, Recent Lamerre paper, \item Slotted e-graphs, nominal logic

\item Atomic completion modulo constructors and primitives. This is a subcase of ground completion where only atomic entities are unioned and with built in injectivity rules for constructors. The semantic domain is $ \mathcal{V} = \{ e_i \}  \uplus \mathcal{P}$ . E-class ids may either be uninterpreted or grounded to a concrete entity. Tie breaking of the union operation on the union find has to be adjusted such that the ground entity is always the root. Unioning two compatible constructors can be resolved by syntactic unification. Unioning two incompatible constructors is considered an error because you are asserting an inconsistent impossibility. Constructors and injectivity can be encoded into the e-graph and rules, but this is a form of baking them in. Having the constructor term reified makes occurs checks and bisimilarity compaction for coalgebraic datatypes easier to perform by ordinary term routines. This is related to the ideas of Co-egraphs \cite{coegraph} which add a cofunction/observation table in addition to the enode function table. The identification of e-ids with values suggests a method to support lambda terms by using Closure values as structured e-ids in an intriguing connection with the ethos of Normalization by Evaluation. In NbE, the semantic domain is extended by syntactic forms representing stuck terms. This can be seen as a method to merge the notion of e-ids and primitives.
% Fixed rewrite rules. If the signature of the "constructors" is disjoint from the other function symbols, one can also have the constructors
% rewrite themselves using smart constructors. To combine them, it sems in general one has to go to full knuth bendix ocmpletion

% Lambda terms are a kind of value

\item A Boolean theory of ground clauses $ \lnot a \lor b \lor c$ or ground Horn clauses $b \land c \implies a$ can be canonized by ordered resolution. This can also be extended to ground superposition clauses $ \lnot(x = y) \lor z = x \lor c = true$. This may give a method to encoding contextual equality. One can also similarly consider using a theory of boolean rings.  %a la colored \cite{singher2023colored}
% the clause c :- a,b is subsmued by c :- a 

\item Disjoint signature of terms with a convergent rewrite system. One not entirely satisfying way of baking in a theory of convergent rewrite rules is if they operate over a disjoint signature from the e-nodes of the egraph. Fuller support (at cost) is to use Knuth Bendix completion as suggested below.

\end{itemize}

It is possible to use theory solvers that may not terminate. One may either hope they do, or terminate their completion process early and accept the incomplete canonizers. The e-graph typically has a self-healing character in the face of incomplete canonization.

\begin{itemize}
\item String Rewriting $<e_1, e_2, ... | R_E>$ . The completion of string rewriting in general will not complete because string rewriting is rich enough to encode turning complete computation, and successful completion would give a decision procedure for . This is an important counterexample as string rewriting can be viewed as ground equations + associativity. Because of this, a fully satisfying solution to an e-graph (which represent ground equations) with baked in associativity is unlikely to be possible. Sequence e-graphs. This may be useful for modelling programs which often need sequencing of side effectful instructions like in Denali. Two related theories are non-commutative rings and finitely presented groups, both of have a flavor of sequences and are not decidable. \item The author is uncertain if it is in the literature if the word problem of differential rings is undecidable, but it probably is. This suggests the possibility however of a method to bake in some notions of calculus and differentiation into the e-graph.

\item Knuth Bendix Completion - One may choose e-ids to be so structured as to be terms themselves.
 Unorientable derived equations can be lifted to be equality saturation rules. This is related to ordered completion/rewriting.
 The ideas around knuth bendix completion are how destructive rewriting (demodulation) can be used while retaining completeness in equtional theorem proving systems. A given set of desirable destructive rewrite rules may be completed ahead of time, but adding in new ground equations in general. 
 The term ordering used to orient and complete the original system can also be determined ahead of time. Strongly term decreasing simplification rules may prefer suing Knuth Bendix orderings which orient roughly speaking by term size. Knuth Bendix completion is such a powerful method that it isn't clear there is an advantage to having the outer e-graph layer.
% Likewise Knuth bendix completion 

\end{itemize}


%An e-graph is a bipartite graph of e-classes and e-nodes. When performing design space exploration on a structure like this, it can be difficult to know when to assign

An e-graph rewriting system may want to mix and match these possibilities. The most straightforward design is to associate a theory with each sort. The sort object is responsible for canonization and maintaining to normalizer data structures. It is not always the case that it makes sense to have one privileged notion of AC for example per sort.

\section{Alternative Approaches}

The above approach is not the only one that might reaosnably go under the moniker e-graphs modulo theories.

\subsection{Extract, Simplify, and Insert}

While the e-class indirection makes theory normalization inside the e-graph confusing, it is known how one would normalize an AST-like term. 
Associativity and commutativity can be dealt with by flattening and sorting an AST into a vector and there are plethora of techniques for beta normalization of terms.

A simple idea is to replace equality saturation rules encoding this behavior with an extraction, simplification, and reinsertion of terms. This allows many directed normalizing steps to be taken at once without bloating the e-graph with administrative junk.

A downside of this technique is that the intermediate states may not really be junk. For AC, the intermediate states may bubble together or collect up terms in such a way as to unlock new optimizations.

Koehler et al describe an "extraction-based substitution" \cite{koehler2022sketchguidedequalitysaturationscaling} method for dealing with substitution of lambda terms. This method works for all rewrite theories for which one wants to avoid junk, not just lambda terms.

\subsection{Constraints via sidecar SMT}

A prominent useful class of automated reasoning solvers are Satisfiability Modulo Theories (SMT) solvers \cite{Barrett2018}. A SAT solver and an e-graph are communication fabric by which theory specific solvers talk to each other. The SAT solver leads to SMT having a backtracking flavor whereas e-graph rewriting has a saturation flavor. A reasonable approach to E-graphs modulo theories can be found in the motivating schematic equation $EMT = SMT - SAT$. If one only asserts unit clauses to an SMT solver then the SAT solver is not needed.

The core technology of SMT solvers is largely the same, but optimized for different use patterns and the first order model semantics is not the desired one for the purposes of optimization.

Upon satisfiability, SMT returns an arbitrary model. Equality saturation specifically wants a justified or minimal model, one in which every equality is forced by the constraints. As an example, given the problem `a = b` `c = d` an SMT solver may return a universe containing only a single element `{a : v!0, b : v!0, c: v!0, d ; d!0}, identifying everything with everything.

The e-matching d the e-matching algorithms are syntactic and furthermore CVC5 and Z3 do not necessarily make theory specific operators like integer addition available to the e-matcher.
Alt-ergo is an exception.

Nevertheless, a technique for EMT is to maintain an equality saturating e-graph and to mirror all the equalities and theory specific facts into a side car SMT solver. The utility of this is to enable SMT queries, specifically in guards of equality saturation rules. Guards are distinguished from multipatterns in that no new pattern variables may be bound in them.


SMT solvers supply a black box for checking equality of two terms modulo the asserted axioms and theories. If one asserts $t \ne s$ and receives an unsat result, then the axioms justify $t = s$ in the minimal model. Given a term bank and bottom-up e-matching, this is sufficient to pattern match, albeit while requiring at worst case $T$ SMT queries. With the addition of using unsat cores, this can be made more query efficient by querying many equalities at once.

\ref{extz3graph} 

% \begin{lstlisting}[language=Lisp]
% (push)
% (assert (not guard))
% (check-sat)
% (pop)
% \end{lstlisting}

This technique has analogs in the form of Constraint Logic Programming \cite{JAFFAR1994503} or Constrained Horn Clauses, or constrained rewriting  \cite{constrainedrewrite}.


\section{Containers and Generalized ENodes}

An alternative to consider from structured e-ids is generalized e-nodes. The basic idea is to merge the concepts of e-node and container primitive.

A standard e-node is a function symbol and a fixed number of ordered children.

Similar to the datatype for terms modulo theories, it is natural to consider an e-node that holds it's children in other data structures like sets, multisets, vectors or polynomials. This is a every so slightly shifted design from putting those structures in the e-ids. During rebuild, the children containers of e-nodes need to be rebuilt. Bottom-up e-matching still works and for some containers top-down e-matching can still work if the interface for enumerating children is extended to be able to return a list.

\section{Just Use Terms Modulo Theories}

Terms are much simpler than e-graphs. While greedy destructive rewriting may miss oppurtunities, one can instead not destroy old terms but keep them around in a term bank for future ocnsideration. This is effectively the same as e-graph techniques without the compaction and sharing (which was the cause of much trouble anyhow). %Brute force equational search is not entirely unreasonable method

Ordered Rewriting can also handle difficult to orient equations in a different way. Ordered rewriting maintains equations like AC and only uses them at rewrite time if they reduce the term according to a term ordering. This guarantees termination and can be used for example as a framework to sort the children of a commutative node. E-graph rewriting techniques should have a superset of the proving or optimization power of ordered rewriting.

Stochastic equational search is another method with a lot to reccommend it. The direction of application of equational rules can be picked probablistically. Some rewrite rules are irreversible, dropping a pattern matched subterm. With the addition of a conservative term bank, even these irreversible rewrite rules can in principle be reversed. If one is going to use a stochastic profile guide extraction method anyway, there is not a super strong reason to maintain an entire e-graph to get a global extraction.

%For better and worse the stochastic nature has a lot of heuristic tuning parameters.

\section{Related Work}

Other examples of e-graph extensions that are evocative of the structured e-id idea are colored e-graphs \cite{singher2023colored} which tag e-ids with a notion of context and slotted e-graphs which  have some of the flavor of both generalized e-ids and generalized e-nodes. 

There have been a number of extensions to rewriting to term rewriting modulo theories. Chapter 11 of Baader and Nipkow has a survey of some options \cite{traatbook}. Of particular interest is Normalizing rewriting \cite{MARCHE1996253} which has the most similarity to the approach here.

String Knuth Bendix systems appear in computational group theory applications. These systems by construction have built in associativity axioms.

Vampire and the E theorem prover are two powerful automated theorem provers with some built in support for AC symbols.

There is a line of work for congruence modulo AC applied in the alt-ergo SMT solver \cite{conchon12lmcs}

The extreme sharing of the e-graph is both its benefit and its problem. There are other term rewriting methods that can handle difficult to orient rules or baked in ness without going

Relational AC matching \cite{relationalac} uses the ideas of Relational E-Matching \cite{zhang2022relationalematching} to encode AC matching into database queries.% is part of a general observation that it is possible to embed different forms of graph matching into relational database queries.  Matching over a term is embeddable into graph matching. The common definition of term nodes has ordered children, whereas many common graph definitions have unordered children.

\section{Conclusion}

Bottoms up! Go forth and saturate!

% Bottoms Up!


\begin{acks}
I'd especially like to thank Cody Roux for many many belligerant grillings about term rewriting. I'd also like to thank Max Willsey, Yihong Zhang, Yisu Remy Wang, Graham Leach-Krouse, Greg Sullivan, Caleb Heibling, Sam Lasser, and Ryan Wisnesky.
\end{acks}

\bibliographystyle{ACM-Reference-Format}
\bibliography{refs}



\appendix

\section{Appendix 1}

% \section{Containers and Generalized ENodes}

% E-graphs systems are the interplay of two entities, enodes and eclasses. A natural approach to generalize e-graphs is to consider generalizing one or both.

% Egg extended the base E-graph with attached lattice values, Egglog has extended with the notion of primtitive and containers.


% Many of the design ideas and issues appear also in the context of pure terms or hash conses. Enodes are basically the same thing a nodes of a term. Eids are to same thing as the unique identifiers of a hashcons, just with extra equality capability.

% Enodes are the data of a model corresponding to function symbols and the eids are the values of the model. The standard simple eid is an integer of some kind, but the fact they are integers are irrelevant. They are uninterpreted entities whose only property is that they can be compared for equality.

% Off by One design problems. Edges and vertices. Material flux is most naturally associated with the faces of cells, not the volumes or vertices. One can mush the different pieces around, wherever there is an enode there is an eclass nearby. A goal is to get the right conceptual picture and the hopefully things will lock into place. This is conceptual advance, and kind of is the only way to advance from our ancestors.

% The most standard notion of enodes is a function symbol combined with a fixed number of ordered children.

% Enodes can be generalize such that their children eids are held in different data structures.
% If enodes hold their children in vectors this supports Associativity to some degree, multisets support AC, and sets support ACI.

% In the egg rust library, a generic notion of enode is enabled by the use of the Language trait. This trait requires the constructors of possible enodes to be distinguishable, teh children enumerable. The capaibiulities are necessary to canpnized and ematch through the enode.

% The \lstinline{fn children(&self) -> &[Id];} This can be generalized into \lstinline{fn children(&self) -> Vec<Vec<Id>>;}

% Some e-graph systems support primitive types such as 64 bit integers or concrete strings. Another primitive type that can be supported is Vectors


\subsection{A No Go Counterexample}

String rewriting is finding pattern strings and replacing them. It is fairly straightforward to encode a turing machine tape and head into a string \cite{traatbook}.
The transitions of the head can be encoded as string rewrite rules. Because of this, string rewriting is Turing complete and we can't expect all
equational theories generated by string equations to be decidable.

%The typical definition of string rewriting does not include a notion of variable.

String rewriting can be shallowly embedded into term rewriting in multiple ways. One method uses a binary concatenation  operation that is associative $(X \cdot Y) \cdot Z = X \cdot (Y \cdot Z)$. The remaining equations of a particular string rewrite system can be represented as ground term equations. If we take the view that an E-graph represents decision procedure for ground term equations then even before discussions of e-matching, an e-graph with a baked in notion of associativity cannot exist.

No No-Go theorem is ever actually a No-Go. It just becomes clear that you need to side step the preconditions or framing of the theorem to achieve what you want. Undecidability of the halting problem does disprove the existence of program analysis and the failure of single layer neural networks to encode xor does not say much about the ability of more complex ones.

It is known that \cite{MARCHE1996253} that ground completion modulo theories is undecidable for

\begin{itemize}
\item A
\item Groups
\item ACD
\end{itemize}

\subsection{A Go Go Example}

The meaning of the notation of conventional mathematics is often both contextual and ambiguous. Like English, this is not necessarily an impediment to effective communication between those with trained eyes. The polynomial expression $x^2 + 2x + 1 = 0$ may be making a statement about a particular point $x$, or about a surface, or describing an algebraic fact about elements of any element of some ring $x$. The former is more analogous to a ground equation in therm term rewriting sense and the latter is more like a non-ground universally quantified equation.

A system of polynomial equations can be converted into one that is equivalent but normalizing under the operation of multivariate polynomial division. This sets of polynomials is called a Gr\"{o}bner basis \cite{Cox2015}. The algorithm to do so is Buchberger's Algorithm.
A Gr\"{o}bner basis is the direct analog of a normalizing rewrite system and Buchberger's algorithm is the direct analog of Knuth-Bendix completion.
Like ground Knuth Bendix completion, Buchberger's algorithm is guaranteed to terminate. Unlike ground knuth bendix completion, the complexity is unlikely to be polynomial.

While it is annoying to make the connection truly rigorous, it is clear that Buchberger's algorithm is performing ground completion modulo a baked in theory of rings.

Two special subcases of Buchberger's algorithm occur for linear polynomials and for polynomials consisting of only two monomials.
The first is a perspective on Gaussian elimination. The latter can be seen as ground multiset rewriting or as integer lattice rewriting. It has connections to the theory of integer programming \cite{sturmfels1996grobner}

It is known that ground completion modulo theories is decidable for the theories of \cite{MARCHE1996253}

\begin{itemize}
\item AC
\item ACU
\item ACI
\item ACN
\item ACUI
\item ACUN
\item Boolean Rings
\item Commutative Rings
\item Finite Fields
\item Abelian Groups
\end{itemize}


\subsection{Completeness: Should One Care?}

Naive equality saturation is not a complete equational theorem proving method because it is defeated by equations that have different variables on each side. \cite{remybirkhoff}

Consider the theory $\forall X, foo(X) = a$  $\forall Y, foo(Y) = b$. and the goal $a =^? b$. 
In naive equality saturation, these two equations can only be used as rewrite rules from left to right. Inserting the goal terms $a$ and $b$ into the e-graph's termbank will not be succiient to trigger the rules.

Either using unification in the form of superposition of unordered ocmpletion, or hvaing an enumeration method of ground terms is needed. \cite{GeMoura2009}

The term banks of the EMT methods do not reduce the total number of e-classes if one wants to find the same number of solutions as naive equality saturation. The number of enodes is reduced.

In AC, if one genuinely wants to be able to search over all subsets of a multiset, 


It is the author's belief that with a suffiicient term enumeratiom mechanism, the EMT methods above are also complete. Even the undecidable theories remain goal-complete if one allows for fair scheduling. These are complex claims that it is outside the scope of this paper to actually prove.

The main interest in understand the completeness of an EMT method relative to the naive algorithm not because the naive algorithm has good theoretical properties. it is that naive equality saturation is simple to explain, pragmatic, and efficient. Alternate EMT methods may have greater or less reasoning power than the naive algorithm on different problems. If the method is equally simple and more efficient that is fine.

Completeness is useful if you want absolute guanrantees you have the minimum possible term or that if the saturation saturates there is no solution.

Soundness however is table stakes.

\end{document}
\endinput

